% ==========================================================================
% GENERATED by scripts/split_paper.py. Do not edit this file.
%
% Source: imported_papers/probe-norm/content/body.tex, content/supplement.tex
% Part 7: Synthetic Controls
%
% Edits belong upstream, in the published paper. To change what the
% chapter includes, change imports/bounded.toml and re-run
% the splitter.
% ==========================================================================
\section{Synthetic Controls}
\label{app:synthetic}
To test whether re-projection recovers known errors, we build latents whose state content we control exactly. The states $\mathbf y\in\mathbb R^3$ follow a stable rotating linear system driven by random actions for $40$ frames. We use $300$ trajectories to fit the probes and $100$ each for validation and scoring.

\emph{State in the direction.} We append a constant one to the state, normalize the result to unit length, and embed it with a fixed matrix $Q$ whose columns are orthonormal, giving $\z=cQ[\mathbf y,1]/\|[\mathbf y,1]\|\in\mathbb R^{16}$. The direction of $\z$ determines $\mathbf y$, because the appended constant keeps the scale that normalization would otherwise remove. The scalar $c$ is the length of $\z$ and plays the role of norm drift: changing it moves the latent along its direction without changing the state it encodes. We fit the probes at $c=1$, as on encodings, and test them at other values of $c$ with the same states, as on drifting predictions. We compare the ReLU probe (Section~\ref{app:impl}), ridge regression, a normalized probe, and the reference mean.

\emph{State in the norm.} The second representation, $\z=c\mathbf u_0$, points in the same fixed direction $\mathbf u_0$ for every sample, and its length is the target state, $y=c$. Here the norm is the only carrier of state. A probe that sees only the direction receives the same input for every sample and cannot recover the state.

%%BEGIN_SYNTHETIC
Both ordinary probes fit the reference encodings well ($R^2=$ $0.9998$ for ReLU and $0.9786$ for ridge), yet they respond strongly to the norm (Figure~\ref{fig:suite}A). Multiplying the norm by ten, with the states unchanged, increases the ReLU probe's error $1.0\times10^{5}$-fold and raises ridge error to $78$ times that of the reference mean. The normalized probe is unaffected. Both ordinary probes still detect deliberate changes of the state-carrying direction, so their response to the norm comes on top of their response to the state. In the norm representation, ReLU and ridge recover the state almost perfectly, while the normalized probe, which sees only the direction, scores a held-out $R^2$ of about zero. Re-projection loses the state in the same way, and the radius check detects this case.

\subsection{A ranking check with known state errors}
\label{app:synthetic_ranking}
We vary state error and norm independently in the directional representation: six state perturbation sizes $\epsilon$ and five norms from a tenth to ten times the reference radius. True state error increases with $\epsilon$ and is independent of norm. Among the $375$ pairs with different perturbation sizes per seed, the ordinary ReLU probe orders only $230$--$232$ correctly, depending on the seed. Re-projection to unit radius, the normalized probe, and the neighbor probe recover every ordering. The neighbor probe has held-out $R^2=$ $0.9962$.

We repeat the ranking test at norms of $0.8$, $1$, $1.25$, and $2$ times the reference radius. Re-projected and normalized probes recover every ordering. Adding a state-independent component orthogonal to the encoding span reduces their accuracy to about $90$\% across $3840$ pairs per seed. A control that preserves the in-span norm gives similar accuracy, so the remaining errors arise from how the probes interpret the off-span component. The neighbor probe remains correct by construction because an orthogonal component cannot change the ordering of reference directions.

\subsection{The protocol's decisions}
\label{app:protocol_test}
We apply the protocol's decision rules to five kinds of prediction whose right reading is known, each at an easy and a hard setting and over five seeds. Drifted predictions carry the state in their direction at two to four times the reference radius, so the re-projected error is the right one. Stable predictions keep the reference radius, so the raw error is. In the third kind, part of the state lies in the norm, which re-projection would remove, so the raw error is right again. Uninformative predictions are other frames' encodings at a drifted norm and should be credited with neither error, and predictions with an added component outside the encodings' span make the probe extrapolate, so their reading is inconclusive. The rules keep the raw error when the radius gain exceeds $0.05$ or the factor exceeds one, call the reading inconclusive when fewer than half of the predictions are supported, and credit the re-projected error when the factor lies below $0.9$, provided it beats the re-projected random directions and the reference mean. They reach the right reading in $50$ of $50$ cases (Table~\ref{tab:protocol_test}). For the predictions off the span, the support check decides. Their re-projected error beats the random directions and the reference mean in $10$ of $10$ cases, while the support check finds $0$\% of them supported.

These cases lie far from the thresholds. Along graded settings that cross them (Table~\ref{tab:protocol_graded}), the decision turns to the re-projected error from a drift of $1.25$ times the reference radius, where the factor first falls below $0.9$, and the re-projected error then stays within $3$\% of the known state error. The support check flags an off-span component from $0.10$ of the in-span norm, where the re-projected error exceeds the known state error by only $13$\%, so the check flags the component before it distorts the error much. The re-projected error is credited in every seed once $75$\% of the predictions encode their own frame; at one half, it lies at the reference mean and the seeds split.

\input{tables/table_protocol_graded.tex}

\input{tables/table_protocol_test.tex}

\paragraph{Rankings.} Step 4 compares pairs of models with several seeds each, so we test it on pairs of synthetic models in five synthetic worlds. A model predicts the state five steps ahead from a window's start state and actions, with a known state error, at a norm drifted from the reference radius, and each model has ten seeds. Its null rollouts predict with another window's actions, from another window's start, and from the start alone, as a module near the identity would at initialization. In informative pairs the more accurate model drifts further. Equally accurate pairs share the state error at different norms, uninformative pairs predict other windows' futures, and pairs off the span add a small component outside the encodings' span. The ordinary probe reverses the true ordering of the informative pairs in $5$ of $5$ worlds and resolves the equally accurate pairs in $5$. The protocol confirms the true ordering of the informative pairs in $5$ worlds and reaches the right verdict on $20$ of $20$ pairs (Table~\ref{tab:protocol_ranking}).

\input{tables/table_protocol_ranking.tex}
%%END_SYNTHETIC

%%BEGIN_SUPPLEMENT
